% Propietario conservado: /Users/ruben/Documents/New project/output/EXTREMA_MEDIA_RAZON_RECIPROCIDAD_ES_EN_20260918/ARTICULO/ES/source/nuclear/12.tex
\chapter{Acciones físicas y espectro de la memoria reducida}
\label{x_reciprocidad:nuclear:12}
\section{Genealogía de las secciones dimensionales}

Las coordenadas \(\pi,\varphi,e,\alpha\) son las realizaciones de las salidas correlacionadas construidas en las secciones~\ref{x_reciprocidad:sec:alpha} y~\ref{x_reciprocidad:sec:electron}. En esta sección se utiliza \(\alpha=\alpha_A\), con la orientación y la periodicización de la definición~\ref{x_reciprocidad:def:alpha-operativa}:



\begin{equation}
\alpha=\pi+e-\varphi-4-\kappa_{\rm per},\qquad
\kappa_{\rm per}=\frac{N_K}{1000^{12}-1},
\quad N_K=234543140729659824621058914794146601.
\label{x_reciprocidad:nuc12:eq:2.1}
\end{equation}



La identidad~\eqref{x_reciprocidad:eq:alpha-identidad-completa} se deduce de la recurrencia con acarreo compatible y de su telescopía, demostradas en la proposición~\ref{x_reciprocidad:prop:alpha-telescopia}. El período doce corresponde al registro: una ventana de doce bloques con condición de frontera nula y una sección infinita compatible son objetos diferentes. La igualdad \eqref{x_reciprocidad:nuc12:eq:2.1} utiliza esta última. El registro ordenado \(K\) conserva sus incidencias además de admitir la lectura escalar \(\kappa_{\rm per}\).

El lector de acción, construido en la sección~\ref{x_reciprocidad:iv:action:section}, determina



\begin{equation}
\begin{aligned}
H_5&=\tfrac12\sqrt{1000\alpha/\varphi}-\alpha+
\tfrac9{16}\alpha^2-\tfrac59\alpha^3+
\tfrac7{48}\alpha^4-\tfrac1{54}\alpha^5,\\
D_A&=\exp(-100\pi\alpha/9),\\
\mathcal C_\pi&=169\alpha^6+\frac{D_A\alpha^7}{1-D_A\alpha},\\
\eta_{\rm ret}&=H_5-\frac{90}{\pi}\mathcal C_\pi,\\
\hbar_{\rm pre}&=H_5\mathcal A_0,\qquad
\hbar_{\rm ret}=\eta_{\rm ret}\mathcal A_0,\qquad
R_{\rm act}=H_5/\eta_{\rm ret},
\end{aligned}
\label{x_reciprocidad:nuc12:eq:2.2}
\end{equation}



Aquí \(\mathcal A_0=10^{-34}\mathcal U_S\). La década procede del lector determinantal de \(\varphi\alpha^{16}\), cuya construcción y acotación figuran en el teorema~\ref{x_reciprocidad:iv:action:decada}; \(H_5\) y \(\eta_{\rm ret}\) son las mantisas de las dos secciones positivas de acción reducida. En cada sección \(a\in\{\mathrm{pre},\mathrm{ret}\}\) se cumple \(h_a=2\pi\hbar_a\).

La construcción del electrón central, desarrollada en la sección~\ref{x_reciprocidad:sec:electron}, determina la coordenada



\begin{equation}
\begin{aligned}
\Delta_4&=\frac{\pi-e\log\pi}{270}\left(1+\frac{\pi}{729}\right),\\
B_e&=\frac{\sqrt3}{4}
\left[\exp(\varphi/\pi^2)-22\alpha^3\right](1+15\Delta_4),\\
\kappa_e&=80-54\alpha+6\Delta_4,\qquad
\varepsilon_e^\beta=B_eR_{\rm act}^{\kappa_e}.
\end{aligned}
\label{x_reciprocidad:nuc12:eq:2.3}
\end{equation}



El prefactor \(\sqrt3/4\) procede de la norma del conmutador de los proyectores de suma y producto sobre las fibras especificadas en \eqref{x_reciprocidad:eq:el-prefactor}. La ruta central y sus registros determinan los coeficientes de retorno de \eqref{x_reciprocidad:eq:el-kappa} y~\eqref{x_reciprocidad:eq:el-beta}. La realización dimensional es \(E_e^\beta=\varepsilon_e^\beta\mathcal U_E\), \(m_e^\beta=E_e^\beta/c^2\), con la carta energética de \eqref{x_reciprocidad:eq:el-dimensional}.

El reloj dodecafásico determina la base energética de \eqref{x_reciprocidad:eq:el-cuanto-reloj}:



\[
\omega_{108}=\frac{2\pi}{108t_0},\quad
E_{\rm clk}=\hbar_{\rm ret}\omega_{108},\quad
b_e=B_e\frac{\mathcal U_E}{E_{\rm clk}}.
\]



La sustitución de \eqref{x_reciprocidad:nuc12:eq:2.3}, conservando la misma energía física, demuestra



\begin{equation}
\boxed{\frac{m_e^\beta c^2}{\hbar_{\rm ret}}
=b_eR_{\rm act}^{\kappa_e}\omega_{108},\qquad
\frac{m_e^\beta c^2}{\hbar_{\rm pre}}
=b_eR_{\rm act}^{\kappa_e-1}\omega_{108}.}
\label{x_reciprocidad:nuc12:eq:2.4}
\end{equation}



La división por \(h_a\), en lugar de \(\hbar_a\), convierte la frecuencia angular en frecuencia cíclica. La conversión \(\mathcal U_E/E_{\rm clk}\) permanece explícita: una coordenada expresada en \(\mathrm{MeV}\) cambia al sustituir su base energética por el cuanto del reloj.

\section{Normalización de la acción Maxwell--Dirac}

La respuesta constitutiva de la sección~\ref{x_reciprocidad:iii:sec:constitutiva} produce la velocidad \(c\), la impedancia \(Z\), la permitividad \(\varepsilon\) y la permeabilidad \(\mu\), todas positivas, con \(\varepsilon c=Z^{-1}\) y \(\mu c=Z\). La sección positiva de carga de \eqref{x_reciprocidad:iii:eq:charge} satisface



\begin{equation}
e_a^2=\frac{2\alpha h_a}{Z}
=4\pi\alpha\varepsilon\hbar_a c .
\label{x_reciprocidad:nuc12:eq:3.1}
\end{equation}



Aquí \(e_a>0\) designa la carga elemental de la sección \(a\); \(e\) en \eqref{x_reciprocidad:nuc12:eq:2.1}--\eqref{x_reciprocidad:nuc12:eq:2.3} designa la coordenada de Euler. Fijamos una carta homogénea con \(c,Z,\varepsilon,\mu,\hbar_a,e_a\) constantes positivos, la masa \(m=m_e^\beta>0\) y la carga \(q\in\{-e_a,e_a\}\); entonces \(\sigma=q/e_a\in\{-1,1\}\). En la realización de Minkowski de signatura \((+,-,-,-)\), con \(x^0=ct\), la acción Maxwell--Dirac es



\begin{equation}
\mathcal I=\int d^4x\left[
-\frac{F_{\mu\nu}F^{\mu\nu}}{4\mu c}
+\bar\psi\left(i\hbar_a\gamma^\mu
(\partial_\mu+iqA_\mu/\hbar_a)-mc\right)\psi
\right].
\label{x_reciprocidad:nuc12:eq:3.2}
\end{equation}



La componente \(A_0\) es el potencial escalar dividido por \(c\), y \(F=dA\). La representación espinorial se fija mediante las matrices \(g^I\) de \eqref{x_reciprocidad:vii:eq:clifford-material-explicito}, que satisfacen \(\{g^I,g^J\}=2\eta^{IJ}I_4\) para \(\eta=\operatorname{diag}(-1,1,1,1)\). Definimos \(\Gamma^I=-ig^I\), de modo que \(\{\Gamma^I,\Gamma^J\}=-2\eta^{IJ}I_4\), y, en esta carta plana, \(\gamma^\mu=\delta^\mu_I\Gamma^I\), \(\bar\psi=\psi^\dagger\Gamma^0\). Así se fijan conjuntamente la firma, el adjunto y el factor imaginario; las matrices \(\gamma^\mu\) se distinguen del parámetro escalar de Holst utilizado después. Se consideran campos suaves y variaciones de soporte compacto, o condiciones de frontera que anulen los términos de borde. La acción realiza los coeficientes ya generados y la representación material de la sección~\ref{x_reciprocidad:vii:subsec:accion-espinorial-afin}.

Definimos



\[
\ell_a=\frac{\hbar_a}{mc},\quad X=x/\ell_a,\quad
\Psi=\ell_a^{3/2}\psi,\quad
\mathsf a_\mu=\frac{e_a\ell_a}{\hbar_a}A_\mu,\quad
\mathsf f=d\mathsf a.
\]



Entonces



\begin{equation}
\boxed{\frac{\mathcal I}{\hbar_a}=
\int d^4X\left[
-\frac{\mathsf f_{\mu\nu}\mathsf f^{\mu\nu}}{16\pi\alpha}
+\bar\Psi\bigl(i\gamma^\mu(\partial_\mu+i\sigma\mathsf a_\mu)-1\bigr)\Psi
\right].}
\label{x_reciprocidad:nuc12:eq:3.3}
\end{equation}



\begin{proof}
El cambio de volumen aporta \(\ell_a^4\), cada derivada aporta \(\ell_a^{-1}\) y el producto de los campos fermiónicos aporta \(\ell_a^{-3}\). El coeficiente del término de masa es \(mc\ell_a/\hbar_a=1\). El término electromagnético aporta \(\hbar_a/(4\mu c e_a^2)\). Como \(\mu c=Z\) y \(Ze_a^2=4\pi\alpha\hbar_a\), ese coeficiente es \(1/(16\pi\alpha)\). Quedan determinados ambos términos y sus normalizaciones.
\end{proof}

La variación respecto de \(\mathsf a_\nu\), seguida de la sustitución de \eqref{x_reciprocidad:nuc12:eq:2.1}, da



\begin{equation}
\partial_\mu\mathsf f^{\mu\nu}
=4\pi\left(\pi+e-\varphi-4-\kappa_{\rm per}\right)
\sigma\bar\Psi\gamma^\nu\Psi .
\label{x_reciprocidad:nuc12:eq:3.4}
\end{equation}



La cancelación compone las secciones de acción, la construcción electrónica y la respuesta constitutiva; el lector regional y dodecafásico determina el acoplamiento restante. La masa persiste en la longitud \(\ell_a\), y la acción en la conversión dimensional. La normalización de Compton es una realización posterior de esas salidas. En la realización coulombiana del mismo electrón, la longitud de Bohr es \(\ell_a/\alpha\): sustituir \(e_a^2=4\pi\alpha\varepsilon\hbar_a c\) en \(4\pi\varepsilon\hbar_a^2/(m e_a^2)\) da exactamente esa expresión.

La deducción utiliza coeficientes homogéneos fijos. Para coeficientes dependientes del espacio-tiempo, la regla del producto introduce las derivadas de las escalas y del acoplamiento en las ecuaciones transformadas. La variación de una acción con parámetros variables debe conservar esos términos.

\section{Escala areal de la acción gravitatoria y de la entropía}

La realización circular del ciclo de \(108\) pasos, construida en la sección~\ref{x_reciprocidad:v:ant:sec:gravity}, determina



\begin{equation}
t_P=\frac{54}{\pi}t_0,\qquad
G_a=\frac{c^5t_P^2}{\hbar_a},\qquad
\ell_P^2=\frac{\hbar_aG_a}{c^3}=c^2t_P^2.
\label{x_reciprocidad:nuc12:eq:4.1}
\end{equation}



La condición de coincidencia de radios de esa realización es \(G_a m_{108,a}/c^2=R_{108}\), con \(m_{108,a}=\hbar_a\omega_{108}/c^2\). Su solución única fija el factor \((54/\pi)^2\); el radio empleado es \(G_a m/c^2\). La reciprocidad \(G_{\rm pre}/G_{\rm ret}=R_{\rm act}^{-1}\) conserva \(\ell_P^2\) entre ambas secciones.

La sección~\ref{x_reciprocidad:vii:sec:realizacion-cartan} construye la cotetrada y la conexión a partir del transporte ya realizado. En la acción de la sección~\ref{x_reciprocidad:vii:sec:corriente-respuesta-cartan}, fijamos una cotetrada no degenerada \(e^I\), la orientación, la firma \(\eta=\operatorname{diag}(-1,1,1,1)\) y una conexión métrica \(\omega\), con curvatura \(F^{IJ}=d\omega^{IJ}+\omega^I{}_{K}\wedge\omega^{KJ}\). Tomamos \(\gamma\in\mathbb R\setminus\{0\}\) y mantenemos constantes \(\gamma,\Lambda,c,G_a,\hbar_a\) en la variación. Ponemos \(\Sigma^{IJ}=e^I\wedge e^J\) y \(\mathsf P_\gamma=\star+\gamma^{-1}I\); la dualidad interna lorentziana satisface \(\star^2=-I\) sobre bivectores. Las variaciones tienen soporte compacto o condiciones de frontera que anulan el término superficial. La sustitución de \eqref{x_reciprocidad:nuc12:eq:4.1} en \eqref{x_reciprocidad:vii:eq:accion-cartan-holst} da



\begin{equation}
\boxed{\frac{\mathcal I_{\rm grav}}{\hbar_a}
=\frac1{16\pi\ell_P^2}
\left[\int\Sigma_{IJ}\wedge(\mathsf P_\gamma F)^{IJ}
-2\Lambda\int\operatorname{vol}_e\right].}
\label{x_reciprocidad:nuc12:eq:4.2}
\end{equation}



Con la cotetrada y los campos materiales fijos durante la variación de la conexión, sea \(\sigma_{IJ}\) la corriente material definida por
\(\delta_\omega\mathcal I_m=-\tfrac12\int\delta\omega^{IJ}\wedge\sigma_{IJ}\), según \eqref{x_reciprocidad:vii:eq:corriente-variacional}. La misma sustitución en el teorema~\ref{x_reciprocidad:vii:thm:ecuacion-cartan-holst} produce



\begin{equation}
D_\omega(\mathsf P_\gamma\Sigma)
=8\pi\ell_P^2\,\frac{\sigma}{\hbar_a}.
\label{x_reciprocidad:nuc12:eq:4.3}
\end{equation}



\begin{proof}
En la sección \(a\), el coeficiente original es \(1/(2\kappa_a c)\), con \(\kappa_a=8\pi G_a/c^4\). Dividir por \(\hbar_a\) lo convierte en \(1/(16\pi\ell_P^2)\); el término cosmológico conserva su factor dos. En la ecuación de conexión se cumple \(\kappa_a c\hbar_a=8\pi\ell_P^2\). Ambas identidades utilizan la misma sección dimensional.
\end{proof}

El funcional de Barbero de la sección~\ref{x_reciprocidad:v:ant:sec:barbero} conserva su expresión íntegra. Sus canales son \(q_\pm=e^{-x\mp y}\), con \(x=\pi A/180\), \(y=\pi C^*/180\), donde \(A\) y \(C^*\) son las coordenadas angulares expresadas en grados. El dominio \(x>|y|\) asegura \(0<q_\pm<1\):



\begin{equation}
\gamma=\frac{180}{\pi}xy+
12[S_{90}(q_-)-S_{90}(q_+)]
-[S_{120}(q_-)-S_{120}(q_+)],\qquad
S_m(q)=\frac{q^m}{1-q^{3m}}.
\label{x_reciprocidad:nuc12:eq:4.4}
\end{equation}



La incidencia hexádica y octádica fija las multiplicidades \(90\) y \(120\), y la cadena orientada fija los coeficientes \(12\) y \(-1\), como demuestra \eqref{x_reciprocidad:v:ant:eq:gamma}. La construcción areal de la sección~\ref{x_reciprocidad:v:ant:sec:area} da además



\begin{equation}
a_0=\frac{1+2\cos(\pi/18)}3\,\frac y6,\qquad
\widehat{\mathcal A}=\gamma a_0\ell_P^2(N_{90}+N_{120}).
\label{x_reciprocidad:nuc12:eq:4.5}
\end{equation}



La inversión \(y\mapsto-y\) cambia los signos de \(\gamma\) y \(a_0\); su producto es par. El operador \(\mathsf P_\gamma\) conserva, en cambio, su dependencia de la orientación. La invariancia del producto areal se distingue así de la invariancia de cada coeficiente del operador. El espectro areal se realiza en la sección positiva \(\gamma a_0>0\), con los operadores de ocupación \(N_{90},N_{120}\geq0\) construidos en la sección~\ref{x_reciprocidad:v:ant:sec:area}.

Para el horizonte cosmológico esférico de \eqref{x_reciprocidad:vii:eq:horizonte-cosmologico}, de radio \(R>0\) y área \(\mathcal A_H=4\pi R^2\), fijamos una sección \(a\) y abreviamos \(\hbar=\hbar_a\), \(G=G_a\) durante este cálculo y el del rebote que sigue. Entonces



\begin{equation}
\frac{S_H}{k_B}=\frac{\mathcal A_H}{4\ell_P^2},\qquad
T_{\rm cos}=\frac{\hbar c}{2\pi k_BR},\qquad
E_\Lambda=\frac{c^4R}{2G},\qquad T_{\rm cos}S_H=E_\Lambda .
\label{x_reciprocidad:nuc12:eq:4.6}
\end{equation}



La escala areal de \eqref{x_reciprocidad:nuc12:eq:4.2}, \eqref{x_reciprocidad:nuc12:eq:4.3} y \eqref{x_reciprocidad:nuc12:eq:4.6} es la misma. La temperatura corresponde al horizonte cosmológico esférico; la normalización del horizonte de Schwarzschild constituye otra realización. Al variar \(R\) manteniendo \(c,G,\hbar,k_B\) fijos, la diferenciación de \eqref{x_reciprocidad:nuc12:eq:4.6} conserva ambos términos: \(T_{\rm cos}\,dS_H=2\,dE_\Lambda\) y \(S_H\,dT_{\rm cos}=-\,dE_\Lambda\), como en \eqref{x_reciprocidad:vii:eq:diferencial-horizonte}. Se obtiene el diferencial de una identidad entre estados de esa familia, no una sustitución del balance dinámico con otros parámetros variables.

\subsection{Una aplicación dinámica explícita}

El teorema~\ref{x_reciprocidad:vii:thm:amplitud-material-explicita} determina, para un estado espinorial con segundo momento conservado \(q_\varrho>0\), la siguiente amplitud. La conservación se expresa mediante \eqref{x_reciprocidad:vii:eq:conservacion-momento-cuartico}; el corolario~\ref{x_reciprocidad:vii:cor:dominio-positivo-conservado} proporciona una realización concreta en el sector de tres ocupaciones:



\begin{equation}
s_0=\frac{3\pi G\hbar^2}{2c^2V_c^2}
\frac{\gamma^2}{1+\gamma^2}q_\varrho .
\label{x_reciprocidad:nuc12:eq:4.7}
\end{equation}



En la realización espacialmente plana con materia de presión nula y \(\Lambda=0\), la constante \(\rho_0>0\) es una densidad de energía, \(V_c>0\) es el volumen comóvil de referencia y \(E_0=\rho_0V_c>0\). La solución del teorema~\ref{x_reciprocidad:vii:thm:rebote-polvo} es



\[
a_b^3=s_0/\rho_0,\qquad
\tau^2=\frac{s_0c^2}{6\pi G\rho_0^2},\qquad
a(t)=a_b\left(1+\frac{(t-t_b)^2}{\tau^2}\right)^{1/3}.
\]



La restitución de \(c\) en \eqref{x_reciprocidad:vii:eq:dinamica-homogenea} utiliza el coeficiente \(8\pi G/(3c^2)\) cuando la densidad es energética. En particular, \((\dot a/a)^2=(8\pi G/(3c^2))(\rho_0a^{-3}-s_0a^{-6})\). Sustituir \eqref{x_reciprocidad:nuc12:eq:4.7} en los parámetros de la solución demuestra



\begin{equation}
\boxed{\left(\frac{\tau E_0}{\hbar}\right)^2
=\frac{q_\varrho}{4}\frac{\gamma^2}{1+\gamma^2},\qquad
\mathcal V_b=\frac{3\pi}{2}q_\varrho
\frac{\gamma^2}{1+\gamma^2}\ell_P^2\frac{\hbar c}{E_0}.}
\label{x_reciprocidad:nuc12:eq:4.8}
\end{equation}



Aquí \(\mathcal V_b=V_c a_b^3\) es el volumen físico en el rebote. Se cancelan \(G\) y \(V_c\) en la primera igualdad; \(E_0\) permanece como energía material. En el sector conservado \(\Lambda^3\mathbb C^4\), la proposición~\ref{x_reciprocidad:vii:prop:operador-espin-car} demuestra \(\widehat{\mathscr Q}_{\rm sp}=4I\), por lo que \(q_\varrho=4\) para todo estado normalizado soportado en ese sector, y
\(\tau E_0/\hbar=|\gamma|/\sqrt{1+\gamma^2}\).
Se obtiene así una escala dinámica para el estado espinorial y la realización cosmológica especificados.

\section{Prolongación del transporte de memoria: realización cuántica}

La realización cuántica utiliza el acoplamiento reversible de la sección~\ref{x_reciprocidad:nuclear:10}. Consideramos primero un plano canónico real con matriz alternante \(\Omega=\left(\begin{smallmatrix}0&1\\-1&0\end{smallmatrix}\right)\). La partición nonádica determina



\begin{equation}
Q=\frac13\begin{pmatrix}\sqrt8 I&-I\\I&\sqrt8 I\end{pmatrix},\qquad
U_H=Q^{\mathsf T}\operatorname{diag}(I,H)Q
=\frac19\begin{pmatrix}8I+H&\sqrt8(H-I)\\
\sqrt8(H-I)&I+8H\end{pmatrix}.
\label{x_reciprocidad:nuc12:eq:5.1}
\end{equation}



El operador \(H\in\operatorname{Sp}(2,\mathbb R)\) transporta ese plano. La matriz \(Q\) conserva el producto euclídeo y la forma \(\Omega\oplus\Omega\); por ello \(U_H\) es simpléctica. La primera componente define el sistema observado y la segunda el registro de memoria; la transformación completa actúa sobre ambos.

El lazo ordenado construido en la sección~\ref{x_reciprocidad:nuclear:11} es



\begin{equation}
C=\begin{pmatrix}0&-1\\1&-1\end{pmatrix},\quad
P=\begin{pmatrix}1&1\\0&1\end{pmatrix},\quad
H_n=C^{-1}P^{-n}CP^n
=\begin{pmatrix}1+n&n^2\\n&n^2-n+1\end{pmatrix}.
\label{x_reciprocidad:nuc12:eq:5.2}
\end{equation}



Para cada \(n\in\mathbb Z\) se cumple \(\det H_n=1\). En particular, para \(n=1\),


\begin{equation}
H_1=\begin{pmatrix}2&1\\1&1\end{pmatrix}
=F_1^2,\quad F_1=\begin{pmatrix}1&1\\1&0\end{pmatrix}
=P F_{\rm av}P^{-1}.
\label{x_reciprocidad:nuc12:eq:5.3}
\end{equation}



Su espectro es \(\{\varphi^2,\varphi^{-2}\}\), por el reconocimiento posterior del modo áureo ya generado. El entero \(n\) especifica el número orientado de iteraciones en el lazo elegido. La familia de transportes está determinada algebraicamente; su realización como protocolo físico conserva además los datos de preparación y de evolución correspondientes.

\subsection{Estado inicial, transporte y lectura}

La proposición~\ref{x_reciprocidad:prop:ii-covarianza-gaussiana} construye la realización gaussiana de la elipse de acción. En la sección positiva \(\hbar_a>0\), con \(A>0\) y \(|C^*/A|<1\), su covarianza es



\begin{equation}
V_\eta=\frac{\hbar_a}{2}M_\eta M_\eta^{\mathsf T},\qquad
M_\eta=\operatorname{diag}(e^{\eta/2},e^{-\eta/2}),\qquad
\eta=\operatorname{artanh}(C^*/A).
\label{x_reciprocidad:nuc12:eq:5.4}
\end{equation}



En la representación de Schrödinger sobre \(L^2(\mathbb R,dQ)\), con \(\widehat Q\) multiplicación por \(Q\) y \(\widehat P=-i\hbar_a\,d/dQ\), la relación \([\widehat Q,\widehat P]=i\hbar_a I\) vale en el espacio de Schwartz. El estado puro normalizado es



\[
\psi_\eta(Q)=(\pi\hbar_a e^\eta)^{-1/4}
\exp[-Q^2/(2\hbar_a e^\eta)].
\]



La preparación es el producto \(\psi_\eta\otimes\psi_\eta\). El transporte de \eqref{x_reciprocidad:nuc12:eq:5.1} se expresa en la misma carta de covarianza por



\begin{equation}
\widetilde U=(M_\eta\oplus M_\eta)U_H
(M_\eta^{-1}\oplus M_\eta^{-1}).
\label{x_reciprocidad:nuc12:eq:5.5}
\end{equation}



El lema~\ref{x_reciprocidad:app:lem:implementacion-simplectica} proporciona una implementación unitaria de esta matriz simpléctica en dos modos. Las implementaciones que difieren por una fase global producen la misma conjugación del estado. La covarianza se transforma por \(V\mapsto\widetilde U V\widetilde U^{\mathsf T}\); el estado conjunto permanece puro y su traza parcial sobre el segundo modo determina la lectura visible, conforme a la construcción de la sección~\ref{x_reciprocidad:app:gaussianas-espectro}.

\begin{theorem}[Mezcla reducida por deformación relativa]

La covarianza de la lectura visible es



\begin{equation}
\boxed{V_{\rm vis}=
\frac{\hbar_a}{2}M_\eta
\frac{8I+HH^{\mathsf T}}9
M_\eta^{\mathsf T}.}
\label{x_reciprocidad:nuc12:eq:5.6}
\end{equation}



Su valor simpléctico normalizado es



\begin{equation}
\boxed{\nu(H):=\frac{2\sqrt{\det V_{\rm vis}}}{\hbar_a}
=\sqrt{1+\frac8{81}\left[\operatorname{tr}(HH^{\mathsf T})-2\right]}.}
\label{x_reciprocidad:nuc12:eq:5.7}
\end{equation}
\end{theorem}



\begin{proof}
Sean \(T=(8I+H)/9\) y \(D=\sqrt8(H-I)/9\). Los dos modos iniciales son independientes y tienen la misma covarianza. En la carta equilibrada del estado inicial, la covarianza visible normalizada es \(TT^{\mathsf T}+DD^{\mathsf T}\). La expansión da



\[
\begin{aligned}
TT^{\mathsf T}+DD^{\mathsf T}
&=\frac{64I+8(H+H^{\mathsf T})+HH^{\mathsf T}
+8[HH^{\mathsf T}-H-H^{\mathsf T}+I]}{81}\\
&=\frac{8I+HH^{\mathsf T}}9.
\end{aligned}
\]



Esto demuestra \eqref{x_reciprocidad:nuc12:eq:5.6}. Puesto que \(\det M_\eta=1\) y \(\det(HH^{\mathsf T})=1\),
\(\det(8I+HH^{\mathsf T})=65+8\operatorname{tr}(HH^{\mathsf T})\).
La división por \(81\) demuestra \eqref{x_reciprocidad:nuc12:eq:5.7}. Los dos autovalores positivos de \(HH^{\mathsf T}\) son recíprocos; su suma es al menos dos. En consecuencia, \(\nu\geq1\), con igualdad exactamente cuando \(HH^{\mathsf T}=I\).
\end{proof}

La escala \(\hbar_a\) y la deformación \(\eta\) se conservan hasta efectuar el determinante. Su cancelación es covariante porque se transportan conjuntamente estado y operador. La expresión \(HH^{\mathsf T}\) utiliza el producto euclídeo de la carta equilibrada del estado inicial; un cambio de carta transporta también ese producto métrico.

Un transporte \(H\) ortogonal y simpléctico puede tener \(D\neq0\) y, a la vez, conservar el producto gaussiano isotrópico, con \(\nu=1\). El complemento operatorio y el entrelazamiento del estado tienen, por tanto, criterios distintos. En esta preparación, la deformación relativa no isométrica caracteriza exactamente el caso \(\nu>1\).

\subsection{Evaluación exacta del lazo áureo}

Para \eqref{x_reciprocidad:nuc12:eq:5.2},



\begin{equation}
\operatorname{tr}(H_nH_n^{\mathsf T})-2
=n^2(2n^2-2n+5),\qquad
\nu_n^2=1+\frac8{81}n^2(2n^2-2n+5).
\label{x_reciprocidad:nuc12:eq:5.8}
\end{equation}



La identidad resulta de sumar los cuadrados de las cuatro entradas. Para \(n=1\), la suma es siete y



\begin{equation}
\boxed{\nu_1=\frac{11}{9},\quad
\operatorname{Tr}\rho_{\rm vis}^2=\frac9{11},\quad
\bar n_{\rm W}=\frac19,\quad
\frac{S_{\rm vis}}{k_B}=\frac{10}{9}\log10-\log9.}
\label{x_reciprocidad:nuc12:eq:5.9}
\end{equation}



El teorema~\ref{x_reciprocidad:app:thm:espectro-gaussiano} demuestra que el estado gaussiano centrado de un modo con valor simpléctico \(\nu\) es unitariamente equivalente a la densidad geométrica de \eqref{x_reciprocidad:app:eq:estado-geometrico}, cuya ocupación media es \(\bar n_{\rm W}=(\nu-1)/2\). Para \(\nu=11/9\), sus autovalores son



\begin{equation}
\lambda_j=\frac9{10}\,10^{-j},\qquad j=0,1,2,\ldots.
\label{x_reciprocidad:nuc12:eq:5.10}
\end{equation}



La serie suma uno. Su cuadrado suma
\((9/10)^2/(1-10^{-2})=9/11\). Finalmente,


\[
-\sum_j\lambda_j\log\lambda_j
=-\log(9/10)+\log10\sum_jj\lambda_j
=\frac{10}{9}\log10-\log9.
\]



La reducción gaussiana y su espectro están demostrados en la sección~\ref{x_reciprocidad:app:gaussianas-espectro}; las dos series anteriores evalúan exactamente la pureza y la entropía de esta realización. Como el estado conjunto es puro, \(S_{\rm vis}/k_B\) es también su entropía adimensional de entrelazamiento. La transformación global es unitaria, mientras la restricción al primer modo tiene entropía positiva.

El resultado corresponde a la preparación y al transporte especificados. Su identificación con una medición del vacío requiere la realización experimental de ese protocolo. La entropía de una reducción y el calor transferido se mantienen como magnitudes distintas.

\subsection{Composición con el transductor térmico}

En la representación normal de Williamson se elige una frecuencia \(\omega>0\) y se define el lector energético
\(\widehat H_{\rm W}=\hbar_a\omega(\widehat N+\tfrac12)\). La frecuencia forma parte de esta realización energética; la covarianza determina las poblaciones. Con esta elección, \eqref{x_reciprocidad:nuc12:eq:5.10} es un estado de Gibbs cuyo cociente de poblaciones satisface



\[
e^{-\hbar_a\omega/(k_BT_{\rm W})}=\frac1{10},\qquad
T_{\rm W}=\frac{\hbar_a\omega}{k_B\log10}.
\]



Si se toma la frecuencia del reloj de la misma sección, \(\omega=1/t_P=2\pi/T_{108}\), la identidad del transductor térmico \eqref{x_reciprocidad:v:ant:eq:ii-kb-aut-par-termico},
\(\hbar_a/t_P=k_B\Theta_{{\rm clk},a}\log3\), da



\begin{equation}
\boxed{\frac{T_{\rm W}}{\Theta_{{\rm clk},a}}=\frac{\log3}{\log10}.}
\label{x_reciprocidad:nuc12:eq:5.11}
\end{equation}



La sección de acción y el transductor térmico se cancelan en el cociente. El índice \(a\) de \(\Theta_{{\rm clk},a}\) conserva la compatibilidad con la sección \(\hbar_a\), como en \eqref{x_reciprocidad:v:ant:eq:ii-kb-aut-orientaciones}. La temperatura \(T_{\rm W}\) es la equivalente de Gibbs para el Hamiltoniano declarado; la evolución unitaria aislada conserva su carácter reversible. Esta identificación espectral determina un cociente térmico y no postula una dinámica de termalización.

\section{Extensión a dimensión canónica finita arbitraria}

Sea \(H\in\operatorname{Sp}(2d,\mathbb R)\), con \(d\geq1\) finito. Se preparan dos estados gaussianos puros, centrados e idénticos, de \(d\) modos cada uno, con covarianza \(V_0=(\hbar_a/2)MM^{\mathsf T}\), donde \(M\in\operatorname{Sp}(2d,\mathbb R)\). La sección~\ref{x_reciprocidad:app:gaussianas-espectro} especifica esta preparación y el transporte conjunto de la carta por \(M\oplus M\). En la carta equilibrada, el mismo acoplamiento \(U_H\) produce \(W=(8I+HH^{\mathsf T})/9\). El producto \(A=HH^{\mathsf T}\) es positivo, simétrico y simpléctico.

Existe una base ortonormal simpléctica en la que


\[
A=\operatorname{diag}(e^{2r_1},e^{-2r_1},\ldots,
e^{2r_d},e^{-2r_d}),\qquad r_j\ge0.
\]



La identidad \(A\Omega A=\Omega\), con \(\Omega=\bigoplus_{j=1}^d\left(\begin{smallmatrix}0&1\\-1&0\end{smallmatrix}\right)\), implica que \(\Omega\) lleva el autoespacio de \(\lambda\) al de \(\lambda^{-1}\). Se eligen bases ortonormales emparejadas en esos espacios. El autoespacio de \(\lambda=1\) es \(\Omega\)-invariante y admite parejas canónicas ortonormales. Su reunión construye la base requerida.

La matriz \(W\) es diagonal en esa base. Los productos de cada pareja determinan sus valores simplécticos:



\begin{equation}
\boxed{\nu_j^2=
\frac{(8+e^{2r_j})(8+e^{-2r_j})}{81}
=1+\frac{32}{81}\sinh^2r_j.}
\label{x_reciprocidad:nuc12:eq:6.1}
\end{equation}



El teorema~\ref{x_reciprocidad:app:thm:espectro-gaussiano} compone la reducción por traza parcial con esos valores simplécticos:



\begin{equation}
\boxed{\operatorname{Tr}\rho_{\rm vis}^2=\prod_{j=1}^d\nu_j^{-1},
\qquad
\frac{S_{\rm vis}}{k_B}=\sum_{j=1}^d
\left[\frac{\nu_j+1}{2}\log\frac{\nu_j+1}{2}
-\frac{\nu_j-1}{2}\log\frac{\nu_j-1}{2}\right].}
\label{x_reciprocidad:nuc12:eq:6.2}
\end{equation}



Se adopta \(0\log0=0\) por continuidad. La prueba del teorema~\ref{x_reciprocidad:app:thm:espectro-gaussiano} transforma unitariamente el estado en un producto de \(d\) densidades geométricas: la pureza del producto es el producto de las purezas y su entropía es la suma de las entropías. El resultado vale para cada \(d\) finito y cada transporte \(H\) del dominio, manteniendo todas las ocupaciones del espacio de Fock.

En el lazo \eqref{x_reciprocidad:nuc12:eq:5.3}, \(r=2\log\varphi\) y \(\sinh^2r=5/4\), lo que reproduce \eqref{x_reciprocidad:nuc12:eq:5.9}. Así, el espectro del transporte y la partición nonádica \(8+1\) determinan el espectro de la lectura reducida; las secciones de acción fijan su realización dimensional.

La dimensión de una realización excepcional y su transporte proceden del constructor correspondiente. La ecuación \eqref{x_reciprocidad:nuc12:eq:6.1} determina la ley espectral después de realizar sus modos canónicos y la preparación gaussiana especificada.

\subsection{Identidad operatoria entre memoria y estructura compleja}

En la carta canónica equilibrada, sea \(J=\Omega\), con \(J^2=-I\) y \(J^{\mathsf T}=-J\). Esta matriz realiza la estructura compleja canónica del espacio de fases, y \(D=\sqrt8(H-I)/9\) es el bloque de memoria. Para la misma preparación y el mismo transporte de carta de la sección~\ref{x_reciprocidad:app:gaussianas-espectro}, se cumple



\begin{equation}
\boxed{-(JW)^2=I+[D,J][D,J]^{\mathsf T},\qquad
W=\frac{2}{\hbar_a}M^{-1}V_{\rm vis}M^{-\mathsf T}.}
\label{x_reciprocidad:nuc12:eq:6.3}
\end{equation}



\begin{proof}
Las identidades simplécticas \(HJH^{\mathsf T}=J\) y \(H^{\mathsf T}JH=J\), junto con \(A=HH^{\mathsf T}\), dan \(JAJ^{\mathsf T}=A^{-1}\). Si \(C=HJ-JH\), la multiplicación de los cuatro sumandos produce



\[
CC^{\mathsf T}
=A+JAJ^{\mathsf T}
-HJH^{\mathsf T}J^{\mathsf T}
-JHJ^{\mathsf T}H^{\mathsf T}
=A+A^{-1}-2I.
\]



Por otra parte,



\[
-(JW)^2
=\frac{(8I+A^{-1})(8I+A)}{81}
=I+\frac8{81}(A+A^{-1}-2I).
\]



Como \([D,J]=\sqrt8\,C/9\), se obtiene \eqref{x_reciprocidad:nuc12:eq:6.3}. La identidad vale en cada dimensión canónica finita y para cada \(H\) del dominio, incluidos los transportes que mezclan los modos.
\end{proof}

Los autovalores de ese operador son los \(\nu_j^2\), cada uno con multiplicidad dos. La fórmula~\eqref{x_reciprocidad:app:eq:pureza-general} expresa entonces la pureza directamente en términos del bloque de memoria:



\begin{equation}
\boxed{\operatorname{Tr}\rho_{\rm vis}^2
=\det\!\left(I+[D,J][D,J]^{\mathsf T}\right)^{-1/4}.}
\label{x_reciprocidad:nuc12:eq:6.4}
\end{equation}



En un modo,



\begin{equation}
\nu^2=1+\frac12\|[D,J]\|_F^2.
\label{x_reciprocidad:nuc12:eq:6.5}
\end{equation}



En este protocolo, \(D\) conmuta con \(J\) exactamente cuando la lectura reducida conserva la pureza. Para el lazo \(H_1\), \(\|[D,J]\|_F^2=80/81\), y \eqref{x_reciprocidad:nuc12:eq:6.5} da \(\nu^2=121/81\), de donde resulta la pureza \(9/11\).

El conmutador utiliza la estructura compleja compatible con la covarianza inicial y su producto métrico. Un cambio canónico transporta conjuntamente los tres datos; la transpuesta de \eqref{x_reciprocidad:nuc12:eq:6.3} se refiere a la carta equilibrada. La construcción previa del elemento imaginario en el núcleo formal precede a esta representación, que especifica su actuación canónica.

La evolución conjunta conserva la pureza mediante su implementación unitaria; el conmutador determina la pérdida de pureza de la lectura parcial. La información del estado conjunto permanece en los dos factores y en sus correlaciones. La fórmula determina la pureza de la reducción, que es un funcional no lineal de su operador densidad.

La descomposición real respecto de \(J\) distingue sus componentes compleja lineal y compleja antilineal:



\[
D_{\rm lin}=\frac{D-JDJ}{2},\qquad
D_{\rm anti}=\frac{D+JDJ}{2}.
\]



La primera componente conmuta con \(J\) y la segunda anticonmuta con \(J\). Las identidades
\([D,J]=2D_{\rm anti}J\) y \(JJ^{\mathsf T}=I\) implican



\begin{equation}
\boxed{-(JW)^2=I+4D_{\rm anti}D_{\rm anti}^{\mathsf T},\qquad
\operatorname{Tr}\rho_{\rm vis}^2
=\det(I+4D_{\rm anti}D_{\rm anti}^{\mathsf T})^{-1/4}.}
\label{x_reciprocidad:nuc12:eq:6.6}
\end{equation}



La componente compleja antilineal del bloque de memoria, definida respecto de la estructura compleja del estado inicial, determina exactamente la pureza y el espectro de mezcla de este protocolo. La componente compleja lineal puede transportar memoria conservando la pureza. La distinción se refiere a la linealidad respecto de \(J\); la complejidad combinatoria y la clasificación aritmética de constantes son propiedades de otros objetos.
